\section{Dados}
\label{sec:dados}

\subsection{Fontes}

Uso três fontes de dados. Primeira: os microdados trimestrais da Pesquisa Nacional por Amostra de Domicílios Contínua (PNAD Contínua) do IBGE \citep{ibgepnad}, referentes ao 1º trimestre de 2026 (2026Q1), totalizando mais de 228 mil indivíduos ocupados com pesos amostrais oficiais ($V1028$). Extraio ocupação (COD a 3 dígitos), condição de formalidade/informalidade, rendimento habitual do trabalho principal, anos de escolaridade, idade, sexo, cor/raça e unidade federativa (27 UFs). Segunda: o índice de exposição ocupacional à IA de \citet{eloundou2024gpts}, construído a partir das tarefas detalhadas da taxonomia americana O*NET-SOC. Terceira: os cruzamentos ocupacionais oficiais ISCO-08$\leftrightarrow$SOC publicados pelo Bureau of Labor Statistics \citep{blscrosswalk}, que conectam a COD (compatível com a ISCO-08) ao sistema O*NET-SOC.

\subsection{Construção do cruzamento e do índice}

A exposição $\theta_j \in [0,10]$ de cada ocupação brasileira corresponde à média ponderada dos escores \texttt{dv\_rating\_beta} de \citet{eloundou2024gpts} dos códigos SOC mapeados, calibrada pela estrutura de postos de trabalho observada. O cálculo é estritamente determinístico, reproduzível por rotinas públicas em Python e sem geração probabilística de conteúdo por LLMs. Ocupações não cobertas pelo escopo original (como forças armadas) recebem valor nulo (\texttt{null}) de forma explícita, sem imputações artificiais.

\subsection{Amostra e variáveis}

A amostra econométrica final é composta por $227.629$ indivíduos ocupados em 122 ocupações com escore válido de exposição e informações completas para todas as covariáveis mincerianas. A taxa de cobertura atinge 99{,}7\% do emprego brasileiro observado. As variáveis centrais são:
\begin{itemize}
    \item \textbf{Exposição à IA ($\theta_j$):} índice de 0 a 10 atribuído à ocupação do trabalhador.
    \item \textbf{Escolaridade ($e_i$):} anos completos de estudo do indivíduo ($0$ a $16$).
    \item \textbf{Informalidade ($\text{informal}_i$):} variável indicadora ($1$ se sem carteira assinada, sem CNPJ no conta-própria/empregador ou trabalhador familiar auxiliar; $0$ caso contrário), conforme o critério metodológico oficial do IBGE.
    \item \textbf{Controles mincerianos:} idade, idade ao quadrado, gênero ($1$ para feminino) e variáveis indicadoras para categorias de cor/raça (preta, parda, amarela, indígena; base branca).
    \item \textbf{Formalidade regional:} taxa de formalidade média da UF do trabalhador, calculada pelo método \emph{leave-one-out} (excluindo a observação do próprio trabalhador) para evitar circularidade mecânica.
\end{itemize}

\subsection{Estatísticas descritivas}

A Tabela~\ref{tab:descritivas} resume as características descritivas da amostra de indivíduos ocupados em 2026Q1, ponderadas pelos pesos amostrais da PNAD Contínua. A exposição média é de 2{,}86 (escala 0--10) e 35,2\% da força de trabalho ocupada (35{,}7 milhões de trabalhadores) situa-se em ocupações com $\theta \le 2$. A escolaridade média alcança 11,6 anos e a taxa de informalidade 40,5\%.

\begin{table}[htbp]
\centering
\caption{Estatísticas descritivas da amostra de indivíduos ocupados (PNAD Contínua, 2026Q1; $N = 227.629$). Médias, desvios-padrão, mínimos e máximos ponderados pelos pesos amostrais. Escolaridade em anos completos de estudo; informalidade em porcentagem; rendimento habitual mensal em reais.}
\label{tab:descritivas}
\begin{tabular}{lccccc}
\toprule
Variável & N & Média & DP & Mín & Máx \\
\midrule
% AUTO-GERADO por scripts/build_paper_tables.py — nao editar a mao
Exposi\c c\~ao \`a IA (0--10) & 227.629 & 2,86 & 1,85 & 0,10 & 8,60 \\
Escolaridade (anos) & 227.629 & 11,59 & 3,81 & 0,00 & 16,00 \\
Informalidade (\%) & 227.629 & 40,5 & 49,1 & 0,0 & 100,0 \\
Rendimento habitual (R\$) & 227.629 & 3.556 & 5.443 & 0 & 400.000
 \\
\bottomrule
\end{tabular}
\end{table}

A Tabela~\ref{tab:maiores} detalha as dez maiores ocupações do país, evidenciando o contraste estrutural: enquanto serviços domésticos combinam exposição quase nula ($0{,}1$) e informalidade de 54{,}3\%, ocupações de apoio administrativo registram exposição elevada ($5{,}8$) e baixa informalidade ($18{,}6\%$).

\begin{table}[htbp]
\centering
\caption{As dez maiores ocupações do Brasil (2026Q1): emprego em milhões, exposição à IA, escolaridade média (anos), informalidade (\%) e rendimento médio (R\$).}
\label{tab:maiores}
\small
\resizebox{\textwidth}{!}{%
\begin{tabular}{lccccc}
\toprule
Ocupação & Empregos & Exposição & Anos est. & Informal. & Renda \\
\midrule
% AUTO-GERADO por scripts/build_paper_tables.py — nao editar a mao
Comerciantes e vendedores de lojas & 7,3 & 4,1 & 11,7 & 31,1 & 3.112 \\
Trabalhadores domésticos e outros trabalhadores d... & 6,7 & 0,1 & 9,0 & 54,3 & 1.548 \\
Escriturários gerais & 3,9 & 5,8 & 13,5 & 18,6 & 2.956 \\
Condutores de automóveis, caminhonetes e motocicl... & 3,3 & 2,9 & 11,1 & 67,4 & 2.650 \\
Trabalhadores da construção civil em obras estrut... & 3,2 & 0,9 & 8,3 & 72,2 & 2.668 \\
Outros vendedores & 3,0 & 2,8 & 11,5 & 50,8 & 2.585 \\
Agricultores e trabalhadores qualificados em ativ... & 3,0 & 1,4 & 7,8 & 84,9 & 2.407 \\
Cabeleireiros, especialistas em tratamento de bel... & 2,6 & 2,3 & 11,6 & 69,1 & 2.366 \\
Condutores de caminhões pesados e ônibus & 2,5 & 3,1 & 10,0 & 26,1 & 3.532 \\
Professores do ensino fundamental e pré-escolar & 2,3 & 3,0 & 15,6 & 20,2 & 4.193
 \\
\bottomrule
\end{tabular}%
}
\end{table}

\subsection{Limitações metodológicas e intensidade de capital}

Três limitações merecem explicitação transparente no corpo do texto. Primeira: a exposição é mensurada no nível ocupacional, não refletindo a heterogeneidade intra-ocupacional de tarefas entre trabalhadores com a mesma COD. Segunda: o índice de \citet{eloundou2024gpts} reflete a estrutura tecnológica da economia americana. Em países em desenvolvimento como o Brasil, onde a relação capital-trabalho é substancialmente menor, ocupações formais homônimas podem incorporar mais rotinas manuais ou analógicas de suporte (\emph{task downgrading}), fazendo com que o escore O*NET represente um teto superior para a verdadeira exposição efetiva brasileira. Terceira: os dados capturam equilíbrios transversais de alocação (\emph{sorting}), não identificando impactos causais da difusão tecnológica sobre salários ou transições ocupacionais.
